% ===========================================================================
% Sample lecture — the kit's content subset demonstrated in one short file.
% This is the worked skeleton: copy its shape for lecture01.tex and delete it
% once the course is underway.  Authoring rules: conventions.md and CLAUDE.md.
% ===========================================================================
\kitlecturetitle{The Kit Content Subset, By Example}

\begin{frame}
  \titlepage
  \begin{center}\small\courseurl\end{center}
\end{frame}

\announcementsframe

\section{Structures}

\begin{frame}{Blocks, Lists, And Emphasis}
  \begin{block}{A named block}
    Body text with \alert{alerted emphasis} and inline math $E=\hbar\omega$.
  \end{block}
  \begin{itemize}
    \item<1-> First point, visible from the start.
    \item<2-> Second point, revealed later — the web shows the final state.
  \end{itemize}
\end{frame}

\begin{frame}{Columns And Case Labels}
  \begin{columns}
    \begin{column}{0.55\textwidth}
      Case \case{1}: the gap stays positive.
      Case \case{2}: the gap changes sign.
    \end{column}
    \begin{column}{0.45\textwidth}
      Columns stack top-to-bottom on the web, in reading order.
    \end{column}
  \end{columns}
\end{frame}

\section{Math And Figures}

\begin{frame}{A Derivation, Step By Step}
  Derivations are transcribed as written — every step shown:
  \begin{align*}
    \braket{u_{\vec k}}{u_{\vec k}} &= 1 \\
    \partial_{k}\,\braket{u_{\vec k}}{u_{\vec k}} &= 0 .
  \end{align*}
\end{frame}

\begin{frame}{A Figure Born With Alt Text}
  \fig{figures/lecture00/square-lattice}{lecture00/square-lattice}
  \onslide<2->{The description lives in \texttt{alt/lecture00.tex}, keyed by
  owning lecture and figure name.}
\end{frame}

\section{Takeaways}

\begin{frame}{Takeaways}
  \begin{enumerate}
    \item One content file, two targets, no transform.
    \item Alt text and announcements live in sidecars, not lecture bodies.
  \end{enumerate}
\end{frame}
